% Part: first-order-logic
% Chapter: axiomatic-deduction
% Section: deduction-theorem-quantifiers

\documentclass[../../../include/open-logic-section]{subfiles}

\begin{document}

\olfileid{fol}{axd}{ddq}

\olsection{قضیهٔ استنتاج با سورها}

\begin{thm}[قضیهٔ استنتاج]
\ollabel{thm:deduction-thm-q} اگر $\Gamma \cup \{!A\} \Proves !B$،
آنگاه $\Gamma \Proves !A \lif !B$.
\end{thm}

\begin{proof}
باز بر طول !!{derivation}~$!B$ از $\Gamma \cup \{!A\}$ استقرا
می‌کنیم.

برهان پایهٔ استقرا با پایهٔ برهان
~\olref[ded]{thm:deduction-thm} یکسان است.

برای گام استقرا، باز فرض کنید !!{derivation}~$!B$ از
$\Gamma \cup \{!A\}$ با گام~$!B$ پایان یابد که با قاعدهٔ استنتاجی توجیه
شده است. اگر قاعدهٔ استنتاج وضع مقدم باشد، مانند برهان
\olref[ded]{thm:deduction-thm} پیش می‌رویم. اگر قاعدهٔ استنتاج \QR باشد،
می‌دانیم $!B \ident !C \lif \lforall[x][!D(x)]$ و !!a{formula} به‌شکل
$!C \lif !D(a)$ پیش‌تر در !!{derivation} آمده است، که در آن $a$ در
~$!C$، $!A$ یا $\Gamma$ رخ نمی‌دهد. پس داریم
\begin{align*}
  \Gamma \cup \{!A\} & \Proves !C \lif !D(a),\\
  \intertext{و فرض استقرا اعمال می‌شود؛ یعنی داریم}
    \Gamma & \Proves !A \lif (!C \lif !D(a)).\\
  \intertext{بنا بر}
  & \Proves (!A \lif (!C \lif !D(a))) \lif ((!A \land !C) \lif !D(a))\\
  \intertext{و با وضع مقدم به دست می‌آوریم}
  \Gamma & \Proves (!A \land !C) \lif !D(a).\\
  \intertext{چون شرط متغیر ویژه همچنان برقرار است، می‌توانیم گامی را که با \QR توجیه شده است به این !!{derivation} بیفزاییم و به دست آوریم}
    \Gamma & \Proves (!A \land !C) \lif \lforall[x][!D(x)].\\
    \intertext{همچنین داریم}
    & \Proves ((!A \land !C) \lif \lforall[x][!D(x)]) \lif (!A \lif (!C \lif \lforall[x][!D(x)])),\\
    \intertext{پس با وضع مقدم،}
    \Gamma & \Proves !A \lif (!C \lif \lforall[x][!D(x)]),
\end{align*}
یعنی $\Gamma \Proves !A \lif !B$.

حالتی را که $!B$ با قاعدهٔ \QR توجیه شده اما به‌شکل
$\lexists[x][!D(x)] \lif !C$ است، به‌عنوان تمرین واگذار می‌کنیم.
\end{proof}

\begin{prob}
  برهان \olref[fol][axd][ddq]{thm:deduction-thm-q} را کامل کنید.
\end{prob}

\end{document}
